\chapter{Тягодинамический расчет}

\section{Внешняя скоростная характеристика}

\includeimage
    {Внешняя скоростная характеристика} % Имя файла без расширения (файл должен быть расположен в директории inc/img/)
    {f} % Обтекание (без обтекания)
    {h} % Положение рисунка (см. figure из пакета float)
    {1\textwidth} % Ширина рисунка
    {Внешняя скоростная характеристика 1NZ-FE.} % Подпись рисунка

\section{Определение передаточных чисел трансмиссии}

Целью данного расчёта является определение передаточных чисел трансмиссии автомобиля. Так как диапазон передаточных чисел бесступенчатого вариатора известен, задача расчёта сводится к подбору передаточного числа главной передачи. Для этого задаются максимальная линейная скорость автомобиля $V_{max}$, радиус колеса $R$, а также минимальное передаточное отношение вариатора $i_{min}^{CVT}$. Кроме того, используется характеристика двигателя для выбора оборотов при максимальной мощности.

\subsection*{Расчётные соотношения}

Длина окружности колеса:
\[
l = 2\pi R
\]

Частота вращения колеса при максимальной скорости:
\[
n_{max}^{\text{колесо}} = \frac{V_{max}}{l}
\]

Из массива данных двигателя определяется частота вращения при максимальной мощности:
\[
n_{max}^{P} = n_{\text{дв}}(\max P).
\]

Требуемое передаточное число главной передачи получается из условия равенства частоты вращения двигателя при пиковой мощности и частоты вращения колёс, скорректированной на минимальное передаточное отношение вариатора:
\[
i_{\text{ГП}} = \frac{n_{max}^{P}}{n_{max}^{\text{колесо}} \, i_{min}^{CVT}} = 8.23
\]

\subsection*{Результат}

Подставляя заданные параметры ($V_{max}$, $R$), а также данные двигателя и вариатора, вычисляется искомое передаточное число главной передачи. Диапазон передаточных чисел трансмиссии получается следующим:
\[
i_{\text{трансм}} \in [3.51, 19.27]
\]

\section{Тяговый баланс}

Тяговый баланс представляет собой сравнение тяговых возможностей автомобиля и сил сопротивления движению. В данном расчёте строятся границы тяговых точек вариатора для разных передаточных чисел, чтобы определить диапазон доступной тяги на ведущих колёсах.

\subsection*{Тяговая сила}

Тяговая сила на ведущих колёсах рассчитывается по формуле:
\[
F_{\text{тяги}}(v) = \frac{M_{\text{дв}}(n)\, i_{\text{трансм}}\, \eta_{\text{трансм}}}{r},
\]
где
\begin{itemize}
  \item $M_{\text{дв}}(n)$ — крутящий момент двигателя при частоте вращения $n$,
  \item $i_{\text{трансм}}$ — общее передаточное число трансмиссии (вариатор × главная передача),
  \item $\eta_{\text{трансм}}$ — коэффициент полезного действия трансмиссии,
  \item $r$ — динамический радиус колеса.
\end{itemize}

\subsection*{Силы сопротивления движения}

Суммарная сила сопротивления движению складывается из сопротивления качению и аэродинамического сопротивления:
\[
F_{\text{сопр}}(v) = F_{\text{кат}} + F_{\text{возд}}(v),
\]
где аэродинамическое сопротивление зависит квадратично от скорости:
\[
F_{\text{возд}}(v) = \frac{1}{2} \rho C_x S v^2,
\]
и
\begin{itemize}
  \item $\rho$ — плотность воздуха,
  \item $C_x$ — коэффициент аэродинамического сопротивления,
  \item $S$ — площадь лобовой поверхности автомобиля,
  \item $F_{\text{кат}} = C_r m g$ — сила сопротивления качению.
\end{itemize}

\subsection*{Графическое построение}

Построение тяговой характеристики для бесступенчатых коробок передач немного отличается от аналогичного построения для ступенчатых коробок передач. Основное отличие заключается в том, что тяговдинамическая характеристика во втором случае представляет из себя группу кривых для каждой из передач, а вирпрм соучае - область состояний, ограниченную низшей и высшей передачами бесступенчатой коробки (рис. \ref{img:Тяговая характеристика}).

\includeimage
    {Тяговая характеристика} % Имя файла без расширения (файл должен быть расположен в директории inc/img/)
    {f} % Обтекание (без обтекания)
    {h} % Положение рисунка (см. figure из пакета float)
    {1\textwidth} % Ширина рисунка
    {Тягодинамическая характеристика.} % Подпись рисунка

Для разных передаточных чисел вариатора строятся:
\begin{itemize}
  \item кривая тяги на высшей передаче $i_{min}^{CVT}$,
  \item кривая тяги на низшей передаче $i_{max}^{CVT}$,
  \item изменение тяги на низких оборотах от передаточного числа,
  \item изменение тяги на высоких оборотах от передаточного числа.
\end{itemize}

Кривая сопротивления возрастает с увеличением скорости, а тяговые характеристики убывают с ростом скорости из-за уменьшения доступного крутящего момента на колесе.  

Точка пересечения огибающей тяговых кривых с кривой сопротивления соответствует максимальной возможной скорости автомобиля. Диапазон тяговых возможностей ограничен верхними и нижними кривыми вариатора.

\section{Динамический фактор}

Динамический фактор характеризует способность автомобиля разгоняться и преодолевать сопротивления движению. Он определяется через **свободную тягу** на ведущих колёсах, которая учитывает аэродинамическое сопротивление:
\[
F_{\text{св}}(v) = F_{\text{т}}(v) - F_{\text{возд}}(v),
\]
где
\begin{itemize}
    \item $F_{\text{т}}(v)$ — тяговая сила на колёсах,
    \item $F_{\text{возд}}(v) = \frac{1}{2}\rho C_x S v^2$ — аэродинамическое сопротивление.
\end{itemize}

Динамический фактор определяется как отношение свободной тяги к весу автомобиля:
\[
D(v) = \frac{F_{\text{св}}(v)}{m g}.
\]

Данимическая характеристика представлена на рисунке ниже (рис. \ref{img:Динамическая характеристика}):

\includeimage
    {Динамическая характеристика} % Имя файла без расширения (файл должен быть расположен в директории inc/img/)
    {f} % Обтекание (без обтекания)
    {h} % Положение рисунка (см. figure из пакета float)
    {1\textwidth} % Ширина рисунка
    {Динамическая характеристика.} % Подпись рисунка

\section{Ускорение автомобиля}

Ускорение характеризует способность автомобиля изменять скорость на различных
режимах работы трансмиссии. В расчёте оно определяется непосредственно из
динамического уравнения движения без интегрирования по времени.

Согласно динамическому уравнению, ускорение автомобиля выражается через
свободную тягу:
\[
a(v) = \frac{F_{\text{св}}(v)}{m},
\]
где
\[
F_{\text{св}}(v) = F_{\text{т}}(v) - F_{\text{возд}}(v).
\]

Таким образом, ускорение зависит только от:
\begin{itemize}
    \item тяговой силы на ведущих колёсах $F_{\text{т}}(v)$;
    \item аэродинамического сопротивления $F_{\text{возд}}(v)$;
    \item массы автомобиля $m$.
\end{itemize}

Расчёт выполняется для крайних передаточных чисел вариатора
($i_{\mathrm{CVT}}^{\min}$ и $i_{\mathrm{CVT}}^{\max}$), что позволяет получить
верхнюю и нижнюю границы ускорения в рабочем диапазоне скоростей.

График ускорения представлен на рисунке ниже (рис. \ref{img:Ускорение}):

\includeimage
    {Ускорение} % Имя файла без расширения (файл должен быть расположен в директории inc/img/)
    {f} % Обтекание (без обтекания)
    {h} % Положение рисунка (см. figure из пакета float)
    {1\textwidth} % Ширина рисунка
    {Ускорение автомобиля.} % Подпись рисунка

На графике по оси абсцисс откладывается скорость $v$, а по оси ординат —
ускорение $a(v)$. Область между кривыми для минимального и максимального
передаточных чисел вариатора отражает диапазон возможных разгонных
характеристик автомобиля.

\chapter{Поверочный расчёт прижимной силы конусов вариатора}

Перед выполнением поверочного расчёта необходимо определить, обеспечивает ли максимальное давление в гидросистеме такую силу прижатия конусов, при которой возможно передать требуемый крутящий момент. Для этого вычисляется реализуемая сила сжатия ремня, её нормальная составляющая с учётом угла наклона поверхности, получаемая касательная сила и соответствующий ей момент.

\section*{Дано}
\begin{itemize}
  \item Максимальный передаваемый момент: $M_{\max}=141{,}00\ \text{Н}\cdot\text{м}$.
  \item Радиус точки приложения силы к ремню: $r=0{,}026\ \text{м}$.
  \item Коэффициент сцепления: $\mu=0{,}10$.
  \item Площадь поверхности прижатия: $A=12\,367\ \text{мм}^2 = 0{,}012367\ \text{м}^2$.
  \item Максимальное давление в системе: $p_{\max}=6\ \text{МПа}=6\cdot 10^{6}\ \text{Па}$.
  \item Угол наклона рабочей поверхности: $\alpha=11^\circ$.
\end{itemize}

\section*{Расчёт}

Реализуемая сила сдавливания ремня:
\[
F_{\text{реал}} = p_{\max}\cdot A
= 6\cdot 10^{6}\cdot 0{,}012367
= 74\,202{,}00\ \text{Н}.
\]

Нормальная сила прижатия:
\[
N_{\text{реал}} = F_{\text{реал}}\cdot \cos\alpha
= 74\,202{,}00\cdot \cos 11^\circ
\approx 72\,838{,}70\ \text{Н}.
\]

Толкающая сила ремня:
\[
T_{\text{реал}} = \mu\cdot N_{\text{реал}}
= 0{,}10\cdot 72\,838{,}70
\approx 7\,283{,}87\ \text{Н}.
\]

Реализуемый момент:
\[
M_{\text{реал}} = T_{\text{реал}}\cdot r
= 7\,283{,}87\cdot 0{,}026
\approx 189{,}38\ \text{Н}\cdot\text{м}.
\]

\section*{Вывод}

\[
M_{\text{реал}} \approx 189{,}38\ \text{Н}\cdot\text{м}
\ge
M_{\max}=141{,}00\ \text{Н}\cdot\text{м}.
\]

Максимальное давление $p_{\max}=6\ \text{МПа}$ обеспечивает достаточную прижимную силу для передачи требуемого крутящего момента при угле наклона конусов $\alpha=11^\circ$.

\chapter{Расчет методом конечных элементов}

В случаях, когда аналитические методы расчёта конструкций становятся слишком сложными или невозможными из-за геометрии, нагрузок или условий закрепления, применяется метод конечных элементов (МКЭ). Этот численный метод позволяет разбить сложную конструкцию на множество мелких, простых элементов, для которых можно легко решить уравнения равновесия.

Для расчета была выбран подвижный конус ведомого шкива вариатора (рис. \ref{img:Подвижный конус}). Он одновременно испытывает нагрузки от давления гидросистемы и центробежных сил вращения. 

\includeimage
    {Подвижный конус} % Имя файла без расширения (файл должен быть расположен в директории inc/img/)
    {f} % Обтекание (без обтекания)
    {h} % Положение рисунка (см. figure из пакета float)
    {1\textwidth} % Ширина рисунка
    {Вид на подвижный конус в сечении трехмерных моделей.} % Подпись рисунка

\newpage

Расчет МКЭ выполнен в среде Siemens NX Nastran с использованием тетраэдрических элементов второго порядка. Модель содержит около 60000 элементов, что обеспечивает достаточную точность результатов.

\includeimage
    {Общая сетка разрез} % Имя файла без расширения (файл должен быть расположен в директории inc/img/)
    {f} % Обтекание (без обтекания)
    {h} % Положение рисунка (см. figure из пакета float)
    {0.5\textwidth} % Ширина рисунка
    {Вид на тетрадную сетку моделей симуляции в разрезе.} % Подпись рисунка

Материал модели — сталь AISI 4340, аналог 40ХН2МА, свойства которой заданы в программе:

\begin{table}[h]
    \centering
    \caption{Основные характеристики материала AISI 4340}
    \begin{tabularx}{\textwidth}{|l|X|}
        \hline
        \textbf{Параметр} & \textbf{Значение} \\
        \hline
        Плотность $\rho$ & $7{.}85 \times 10^{-6}$ кг/мм³ \\
        \hline
        Модуль Юнга $E$ & $193\ \text{ГПа}$ \\
        \hline
        Коэффициент Пуассона $\nu$ & $0.284$ \\
        \hline
        Предел текучести $\sigma_y$ & $1178\ \text{МПа}$ \\
        \hline
    \end{tabularx}
\end{table}

\newpage

\section*{Граничные условия и нагрузки}

Для закрепления модели были зафиксированы все степени свободы в области сопряжения с валом. Нагрузки включали давление гидросистемы, распределенное по внутренней поверхности конуса, и центробежные силы, рассчитанные по угловой скорости вращения.

Закрепления были выставлены через ограничение перемещения вала (рис. \ref{img:Ограничение оси}), вместе с которым вращется и по которому перемещается исследуемый конус.

\includeimage
    {Ограничение оси} % Имя файла без расширения (файл должен быть расположен в директории inc/img/)
    {f} % Обтекание (без обтекания)
    {h} % Положение рисунка (см. figure из пакета float)
    {0.5\textwidth} % Ширина рисунка
    {Ограничение перемещений вала.} % Подпись рисунка

\newpage

Также были выставлены заделки на ребрах звеньев ремня (рис. \ref{img:Ограничение ремня}).

\includeimage
    {Ограничение ремня} % Имя файла без расширения (файл должен быть расположен в директории inc/img/)
    {f} % Обтекание (без обтекания)
    {h} % Положение рисунка (см. figure из пакета float)
    {0.5\textwidth} % Ширина рисунка
    {Заделки на гранях звеньяев ремня.} % Подпись рисунка

\newpage

Помимо этого необходимо было обозначить контакты между звеньями и конусом (рис. \ref{img:Контакт ремня и конуса}), шариками и конусом (рис. \ref{img:Контакт шариков и конуса}), а также между шариками и валом (рис. \ref{img:Контакт шариков и оси}) для реализации расчета.

\includeimage
    {Контакт ремня и конуса} % Имя файла без расширения (файл должен быть расположен в директории inc/img/)
    {f} % Обтекание (без обтекания)
    {h} % Положение рисунка (см. figure из пакета float)
    {0.3\textwidth} % Ширина рисунка
    {Обозначение поверхностей контакта звеньев ремня и конуса.} % Подпись рисунка

\includeimage
    {Контакт шариков и конуса} % Имя файла без расширения (файл должен быть расположен в директории inc/img/)
    {f} % Обтекание (без обтекания)
    {h} % Положение рисунка (см. figure из пакета float)
    {0.3\textwidth} % Ширина рисунка
    {Обозначение поверхностей контакта шариков и конуса.} % Подпись рисунка

\includeimage
    {Контакт шариков и оси} % Имя файла без расширения (файл должен быть расположен в директории inc/img/)
    {f} % Обтекание (без обтекания)
    {h} % Положение рисунка (см. figure из пакета float)
    {0.3\textwidth} % Ширина рисунка
    {Обозначение поверхностей шариков и вала.} % Подпись рисунка

\newpage

После задания граничных условий необходимо задаться нагрузками. У нас их две: давление гидросистемы и центробежные силы.

Давление 6 МПа было распределено по всей рабочей поверхности конуса, чтобы реализовать максимальные напряжения в угле конуса (рис. \ref{img:Давление}).

\includeimage
    {Давление} % Имя файла без расширения (файл должен быть расположен в директории inc/img/)
    {f} % Обтекание (без обтекания)
    {h} % Положение рисунка (см. figure из пакета float)
    {0.5\textwidth} % Ширина рисунка
    {Приложение давления нидросистемы к поверхности конуса.} % Подпись рисунка

\newpage

Центробежную силу необходимо задать с помощья вращения. Для этого необходимо вычислить максимальную угловую скорость вращения конуса:

Пиковая частота вращения на ведомом валу (в об/мин):
\[
n_{\max}^{\text{вал}} = n_{\max}^{\text{двигатель}}\cdot i_{\max}
= 6000 \cdot 2.561 = 15366\ \text{об/мин}.
\]

Переводим обороты в градусы в секунду. Связь:
\[
1\ \text{об/мин} = \frac{1}{60}\ \text{об/с} = \frac{360^\circ}{60\ \text{s}} = 6^\circ/\text{s}.
\]
Следовательно:
\[
\omega_{\text{пик}} = n_{\max}^{\text{вал}}\cdot 6^\circ/\text{s}
= 15366 \cdot 6^\circ/\text{s} = 92196^\circ/\text{s}.
\]

Итог:
\[
\omega_{\text{пик}} = 92\,196^\circ/\text{s}
\]

Теперь можно задать центробежную силу через угловую скорость вращения 92 196 градусов в секунду (рис. \ref{img:Вращение}).

\includeimage
    {Вращение} % Имя файла без расширения (файл должен быть расположен в директории inc/img/)
    {f} % Обтекание (без обтекания)
    {h} % Положение рисунка (см. figure из пакета float)
    {0.5\textwidth} % Ширина рисунка
    {Вращение конуса.} % Подпись рисунка

\newpage

По итогу получилась следующая модель с заданными граничными условиями и нагрузками (рис. \ref{img:Симуляция1}) и (рис. \ref{img:Симуляция2}).

\includeimage
    {Симуляция1} % Имя файла без расширения (файл должен быть расположен в директории inc/img/)
    {f} % Обтекание (без обтекания)
    {h} % Положение рисунка (см. figure из пакета float)
    {0.4\textwidth} % Ширина рисунка
    {Вид на модель с приложенными граничными условиями и нагрузками спереди.} % Подпись рисунка

\includeimage
    {Симуляция2} % Имя файла без расширения (файл должен быть расположен в директории inc/img/)
    {f} % Обтекание (без обтекания)
    {h} % Положение рисунка (см. figure из пакета float)
    {0.4\textwidth} % Ширина рисунка
    {Вид на модель с приложенными граничными условиями и нагрузками сзади.} % Подпись рисунка

\newpage

\section*{Результаты расчёта}

По итогу получились следующие эпюры напряжений и деформаций (рис. \ref{img:Деформация}) и (рис. \ref{img:Напряжение}).

\includeimage
    {Деформация} % Имя файла без расширения (файл должен быть расположен в директории inc/img/)
    {f} % Обтекание (без обтекания)
    {h} % Положение рисунка (см. figure из пакета float)
    {0.5\textwidth} % Ширина рисунка
    {Эпюра деформаций, мм.} % Подпись рисунка

\includeimage
    {Напряжение} % Имя файла без расширения (файл должен быть расположен в директории inc/img/)
    {f} % Обтекание (без обтекания)
    {h} % Положение рисунка (см. figure из пакета float)
    {0.5\textwidth} % Ширина рисунка
    {Эпюра напряжений, МПа.} % Подпись рисунка

\newpage

По результатам статического МКЭ-анализа максимальное эквивалентное напряжение в конструкции составило
\[
\sigma_{\max} = 526\ \text{МПа}.
\]

Для выбранного материала AISI~4340 предел текучести равен
\[
\sigma_{\text{y}} = 1178\ \text{МПа}.
\]

Запас прочности определим как отношение предела текучести к максимальному напряжению:
\[
n = \frac{\sigma_{\text{y}}}{\sigma_{\max}} =
\frac{1178}{526} \approx 2.24.
\]

\textbf{Вывод:} конструкция обладает достаточным запасом прочности, так как расчётное значение \( n \approx 2.24 \) превышает единицу и соответствует требованиям к статической прочности.